%!TEX TS-program = pdflatex
%!TEX TS-options = -shell-escape
% % % % %   Die folgenden Zeilen müssen ihre Zeilennummern 4 und 5 behalten !!!    % % % % %
\newcommand{\printpraesenzlsg}{false}
\newcommand{\printloesungen}{false}
\newcommand{\printbewertungen}{false}
% % % % %   \newcommand{\printloesungen}{false}                                    % % % % %
\newcommand{\blattnummer}{1}
%\newcommand{\abgabetermin}{\textcolor{red}{bis 11.04.2022, 09:00 Uhr}}
\documentclass[%
  paper=a4,
  fontsize=10pt,
  ngerman
  ]{scrartcl}

% Basics für Codierung und Sprache
% ===========================================================
  \usepackage{scrtime}
  \usepackage{etex}
  \usepackage{shellesc}
  \usepackage[final]{graphicx}
  \usepackage[utf8]{inputenc}
  \usepackage{babel}
  \usepackage[german=quotes]{csquotes}
  \usepackage{datetime2}
  \usepackage{ifthen}
  \usepackage{dingbat}
% ===========================================================

%% Fonts und Typographie
%% ===========================================================
  \usepackage{charter}
  \usepackage{nimbusmononarrow}
  \usepackage[babel=true,final,tracking=smallcaps]{microtype}
  \DisableLigatures{encoding = T1, family = tt* }
  \usepackage{ellipsis}
% ===========================================================

% Farben
% ===========================================================
  \usepackage[usenames,x11names,final,table]{xcolor}
  \definecolor{fbblau}{HTML}{3078AB}
  \definecolor{blackberry}{rgb}{0.53, 0.0, 0.25}
% ===========================================================

% Mathe-Pakete und -Einstellungen
% ===========================================================
  \usepackage{mathtools}
  \usepackage{amssymb}
  \usepackage[bigdelims]{newtxmath}
  \allowdisplaybreaks
  \usepackage{bm}
  \usepackage{wasysym}
% ===========================================================

% TikZ
% ===========================================================
  \usepackage{tikz}
  \usepackage{tikz-cd}
  \usepackage{pgf}
  \usetikzlibrary{matrix}
  \usetikzlibrary{positioning}
  \usetikzlibrary{arrows}
  \usetikzlibrary{shapes}
  \tikzset{>=Latex}
  \usepackage[mode=image]{standalone}
% ===========================================================

% Seitenlayout, Kopf-/Fußzeile
% ===========================================================
  \usepackage{scrlayer-scrpage}
  \usepackage[top=3cm, bottom=2.5cm, left=2.5cm, right=2cm]{geometry}  
  \pagestyle{scrheadings}
  \clearscrheadfoot 
  \setheadsepline{0.4pt}
  \setfootsepline{0.4pt}
  \setkomafont{pagehead}{\normalfont\footnotesize}
  \setkomafont{pagefoot}{\normalfont\footnotesize}
  \lohead[]{Klausur zur Vorlesung \textit{Mathematik I - Theoretische Grundlagen der Informatik} -- Wintersemester 2022/2023}
  \rohead[]{%
    \ifthenelse{\boolean{loesungen}}{\textbf{Lösungen zu Blatt \blattnummer}}{Blatt \blattnummer}}
  %\rofoot[Version: \DTMnow]{%
  %  \ifthenelse{\boolean{loesungen}}{Version: \DTMnow}{}}
  %\lofoot[\jobname.tex]{%
  %  \ifthenelse{\boolean{loesungen}}{\jobname.tex}{}}
  \cfoot[\thepage]{\thepage}
  \raggedbottom
  \usepackage{setspace}          
  \linespread{1.2}
  \setlength{\parindent}{0pt}
  \setlength{\parskip}{0.5\baselineskip}
  \usepackage[all]{nowidow}
  \usepackage{lscape}
% ===========================================================

% Hyperref
% ===========================================================
  \usepackage[%
    hidelinks,
    pdfpagelabels,
    bookmarksopen=true,
    bookmarksnumbered=true,
    linkcolor=black,
    urlcolor=fbblau,
    plainpages=false,
    pagebackref,
    citecolor=black,
    hypertexnames=true,
    pdfborderstyle={/S/U},
    linkbordercolor=fbblau,
    colorlinks=false,
    backref=false]{hyperref}
  \hypersetup{final}
  \makeatletter
  \g@addto@macro{\UrlBreaks}{\UrlOrds}
  \makeatother
% ===========================================================

% Listen und Tabellen
% ===========================================================
  \usepackage{multicol}
  \usepackage{multirow}
  \usepackage[shortlabels]{enumitem}
  \setlist[enumerate]{font=\bfseries,itemsep=0.25\baselineskip}
  \setlist[itemize]{label=$\triangleright$,itemsep=-0.25\baselineskip}
  \usepackage{tabularx}
  \usepackage{array}
  \usepackage{listings}
  \lstset{breaklines=true,
    basicstyle=\ttfamily,
    frame=tblr,
    language=c,
    numbers=left}
  \usepackage{include/mips}

\definecolor{lightgreen}{rgb}{0,0.7,0}
\definecolor{gray}{rgb}{0.5,0.5,0.5}
\definecolor{mauve}{rgb}{0.58,0,0.82}

\lstdefinestyle{mips}{ %
  language=[mips]Assembler,       % the language of the code
  keywordstyle=\color{blue},          % keyword style
  keywordstyle=[2]\color{magenta},     % keyword style
  keywordstyle=[3]\color{red},          % keyword style
  commentstyle=\color{lightgreen},       % comment style
  stringstyle=\color{mauve},         % string literal style
  escapeinside={\%*}{*)},            % if you want to add a comment within your code
  morekeywords={*,...}               % if you want to add more keywords to the set
}

\newcommand{\mipsinline}[1]{\lstinline[style=mips,columns=fixed]{#1}}
\newcommand{\mipsregister}[1]{\mipsinline{\$#1}}
  \newcommand{\qed}{\hfill \ensuremath{\Box}}
  \newcommand{\cmd}[1]{\texttt{#1}}
% ===========================================================

% Sonstiges
% ===========================================================
  \usepackage[textsize=small,color=Red1!80!OrangeRed1!80]{todonotes}
  \usepackage{lipsum}
% ===========================================================

% Shortcuts and Math Commands
% ===========================================================
  \newcommand{\BB}{\mathbb{B}}      % Bool
  \newcommand{\CC}{\mathbb{C}}      % komplexe Zahlen
  \newcommand{\NN}{\mathbb{N}}      % nat. Zahlen
  \newcommand{\QQ}{\mathbb{Q}}      % rat. Zahlen
  \newcommand{\RR}{\mathbb{R}}      % reelle Zahlen
  \newcommand{\ZZ}{\mathbb{Z}}      % ganze Zahlen
  \newcommand{\oh}{\mathcal{O}}      % Landau-O
  \newcommand{\ol}[1]{\overline{#1}}    % überstreichen
  \newcommand{\wt}[1]{\widetilde{#1}}    % überschlängeln
  \newcommand{\wh}[1]{\widehat{#1}}    % überdachen
  \newcommand{\entspricht}{\mathrel{\widehat{=}}}
  
  \DeclarePairedDelimiter{\absolut}{\lvert}{\rvert}    % Betrag
  \DeclarePairedDelimiter{\ceiling}{\lceil}{\rceil}    % aufrunden
  \DeclarePairedDelimiter{\Floor}{\lfloor}{\rfloor}    % aufrunden
  \DeclarePairedDelimiter{\Norm}{\lVert}{\rVert}      % Norm
  \DeclarePairedDelimiter{\sprod}{\langle}{\rangle}    % spitze Klammern
  \DeclarePairedDelimiter{\enbrace}{(}{)}          % runde Klammern
  \DeclarePairedDelimiter{\benbrace}{\lbrack}{\rbrack}  % eckige Klammern
  \DeclarePairedDelimiter{\penbrace}{\{}{\}}        % geschweifte Klammern
  \newcommand{\Underbrace}[2]{{\underbrace{#1}_{#2}}}   % bessere Unterklammerungen
  
  % Kurzschreibweisen für Faule und Code-Vervollständigung
  \newcommand{\abs}[1]{\absolut*{#1}}
  \newcommand{\ceil}[1]{\ceiling*{#1}}
  \newcommand{\flo}[1]{\Floor*{#1}}
  \newcommand{\no}[1]{\Norm*{#1}}
  \newcommand{\sk}[1]{\sprod*{#1}}
  \newcommand{\enb}[1]{\enbrace*{#1}}
  \newcommand{\penb}[1]{\penbrace*{#1}}
  \newcommand{\benb}[1]{\benbrace*{#1}}
  \newcommand{\stack}[2]{\makebox[1cm][c]{$\stackrel{#1}{#2}$}}

% Übungszettel-Krams
% ===========================================================
  \usepackage[tikz]{mdframed}
  \newmdenv[%
    backgroundcolor = gray!20,
    linewidth=0pt,
    skipabove = 0.5cm,
    skipbelow = 0cm
  ]{notes}
  
  % Lösungsausgabe
  \newboolean{praesenzlsgn}
  \setboolean{praesenzlsgn}{\printpraesenzlsg}
  \newboolean{loesungen}
  \setboolean{loesungen}{\printloesungen}
  % Bewertungsschema-Ausgabe
  \newboolean{bewertungen}
  \setboolean{bewertungen}{\printbewertungen}

  \newcommand{\iforiginal}[1]{\ifthenelse{\boolean{praesenzlsgn}}{}{#1}}
  \newcommand{\ifpraesenzlsg}[1]{\ifthenelse{\boolean{praesenzlsgn}}{#1}{}}
  \newcommand{\ifloesung}[1]{\ifthenelse{\boolean{loesungen}}{#1}{}}
  \newcommand{\ifbewertung}[1]{\ifthenelse{\boolean{bewertungen}}{#1}{}}

  \newcommand{\umbruchoriginal}{\iforiginal{\newpage}}        % Seitenumbruch nur im Original
  \newcommand{\umbruchpraesenzlsg}{\ifpraesenzlsg{\newpage}}  % Seitenumbruch nur in Lösung d. Präsenzaufgabe
  \newcommand{\umbruchlsg}{\ifloesung{\newpage}}              % Seitenumbruch nur in Lösung
  \newcommand{\umbruchbewertung}{\ifbewertung{\newpage}}      % Seitenumbruch nur in Lösung mit Bewertungsschema
  \newcommand{\umbruchnurpraesenz}{\ifthenelse{\boolean{loesungen}}{}{\umbruchpraesenzlsg}}
  \newcommand{\umbruchnurlsg}{\ifthenelse{\boolean{bewertungen}}{}{\umbruchlsg}}

  \usepackage{comment}
  \newenvironment{praesenzlsg}%
    {\ifthenelse{\boolean{praesenzlsgn}}{\textbf{--- Präsenzaufgabe Anfang ---} \mbox{} \newline}{}}%
	{\ifthenelse{\boolean{praesenzlsgn}}{\mbox{} \newline \textbf{--- Präsenzaufgabe Ende ---}}{}}
  \newenvironment{loesung}%
    {\ifthenelse{\boolean{loesungen}}{\textbf{--- Lösung Anfang ---} \mbox{} \newline}{}}%
    {\ifthenelse{\boolean{loesungen}}{\mbox{} \newline \textbf{--- Lösung Ende ---} \ifthenelse{\boolean{bewertungen}}{}{\umbruchlsg}}{}}
  \newenvironment{bewertung}%
    {\ifthenelse{\boolean{bewertungen}}{\textbf{--- Bewertungsschema Anfang ---} \mbox{} \newline}{}}%
    {\ifthenelse{\boolean{bewertungen}}{\mbox{} \newline \textbf{--- Bewertungsschema Ende ---} \umbruchlsg}{}}

  \ifthenelse{\boolean{praesenzlsgn}}{}{\excludecomment{praesenzlsg}}
  \ifthenelse{\boolean{loesungen}}{}{\excludecomment{loesung}}
  \ifthenelse{\boolean{bewertungen}}{}{\excludecomment{bewertung}}

  % Aufgabennummerierung
  \newcounter{aufgabennummer} 
  \setcounter{aufgabennummer}{1}
  \newcommand{\aufgabe}[1]{\textbf{Aufgabe \blattnummer.\theaufgabennummer} \stepcounter{aufgabennummer} \hfill (#1 Punkte)}
  %\newcommand{\aufgabetitel}[2]{\iforiginal{\vspace{0.5cm}} \textbf{Aufgabe \blattnummer.\theaufgabennummer \ (#2)} \stepcounter{aufgabennummer} \hfill (#1 Punkte)}
  \newcommand{\aufgabetitel}[2]{\iforiginal{\vspace{0.5cm}} \textbf{Aufgabe \theaufgabennummer \ (#2)} \stepcounter{aufgabennummer} \hfill (#1 Punkte)}
  \newcommand{\praesenztitel}[1]{\iforiginal{\vspace{0.5cm}} \textbf{Aufgabe \blattnummer.\theaufgabennummer \ (#1)} \stepcounter{aufgabennummer} \hfill \smallpencil}

\newcommand{\danger}{\raisebox{0.5mm}{\fontencoding{U}\fontfamily{futs}\selectfont\char 66\relax}}
\newcommand{\ddanger}[1]{{\danger} \textbf{#1} {\danger}}

% Rechnerstrukturen-Krams
% ===========================================================
  \newcommand{\Band}{\ensuremath{\mathbin{\texttt{AND}}}}
  \newcommand{\Bor}{\ensuremath{\mathbin{\texttt{OR}}}}
  \newcommand{\Bnot}{\ensuremath{\mathbin{\texttt{NOT}}}}
  \newcommand{\Bnand}{\ensuremath{\mathbin{\texttt{NAND}}}}
  \newcommand{\Bnor}{\ensuremath{\mathbin{\texttt{NOR}}}}
  \newcommand{\Bxor}{\ensuremath{\mathbin{\texttt{XOR}}}}
  \newcommand{\Bxnor}{\ensuremath{\mathbin{\texttt{XNOR}}}}
  \newcommand{\Bone}{\ensuremath{\mathbin{\texttt{ONE}}}}
  \newcommand{\Bzero}{\ensuremath{\mathbin{\texttt{ZERO}}}}
  \newcommand{\Binary}[1]{\ensuremath{\texttt{#1}}\xspace}
  \newcommand{\ieee}[3]{\fbox{\ensuremath{\texttt{#1}}}\fbox{\ensuremath{\texttt{#2}}}\fbox{\ensuremath{\texttt{#3}}}}
  \DeclareMathOperator{\DNF}{DNF}
  \DeclareMathOperator{\KNF}{KNF}
  \DeclareMathOperator{\xor}{\oplus}
  \newcommand{\dc}{\star} %don't-care-Wert
 

% Änderungen 2020: Hinweise auf Digitallehre angepasst; Rechnergeschichte entfernt; Blatt 1 und 2 zusammengefasst.
\begin{document}

\iforiginal{\thispagestyle{scrplain}
\vspace*{-3cm}
\begin{minipage}[t][1.1cm][c]{4.5cm}
  \includegraphics[width=4cm]{include/hwr-logo.png}
\end{minipage}
\hfill
\begin{minipage}[t][1.5cm][c]{8cm}
  \begin{center}
  \begin{footnotesize}
    \textsf{Prof. Dr.-Ing. Sebastian Schlesinger} \\[-0.1cm]
    \textsf{Fachbereich 2 - Duales Studium Wirtschaft \& Technik}
    %\url{https://www.hwr-berlin.de}
  \end{footnotesize}
  \end{center}
\end{minipage}
\hfill
%\begin{minipage}[t][1.6cm][c]{3.2cm}
%  \includegraphics[width=2.8cm]{include/ESlogo2}
%\end{minipage}

\vspace*{-0.3cm}
\begin{center} 
  \hrulefill \\[0.1cm]
  {\large Klausur} \\[0.15cm]
  {\huge \bfseries Mathematik I - Theoretische Informatik} \\[0.10cm]
  {HWR Berlin, Wintersemester 2024} \\[-0.4cm]
  \begin{tabular}{lcr}
    \hspace{0.3\textwidth}   & \hspace{0.3\textwidth} & \hspace{0.3\textwidth} \\
    Prof. Dr.-Ing. Sebastian Schlesinger   %&                        & "Ubung zur Klausurvorbereitung    \\ 
    %Blatt \blattnummer     % \abgabetermin          \\ 
  \end{tabular} \\[0.1cm]
  \hrulefill
\end{center}}
\aufgabetitel{$6$}{Mengen und Funktionen}\\
Gegeben seien die Mengen $A=\{1,2,3,\{4\}\}$, $B=\{2,4\}$ und $C=\{\emptyset,\{1\}\}$.
\begin{enumerate}
\item Geben Sie die Menge $A\cup B$ an
\item Geben Sie die Menge $A\cap B$ an
\item Geben Sie die Menge $A\setminus B$ an.
\item Was ist $\mathscr{P}(C)$?
\item Geben Sie eine bijektive Abbildung $f:A\to\mathscr{P}(C)$ an.

\end{enumerate}

\aufgabetitel{$12$}{Relationen}\\
Gegeben seien die Relationen $R,S\subseteq\{1,2,3,4,5,6\}$ mit $R=\{(1,2),(2,3),(4,3),(5,4),(6,4)\}$\\und $S=\{(1,1),(1,2),(2,1)\}$.
\begin{enumerate}[(i)]
    \item Stellen Sie $R$ und $S$ als Graphen und Adjazenzmatrix dar.
    \item Stellen Sie die Relation $R\circ S$ als Graphen dar.
    \item Stellen Sie die Relation $S\circ R$ als Graphen dar.
    \item $R$ ist nicht transitiv. Begründen Sie warum (Gegenbeispiel).
    \item Ergänzen Sie den Graphen von $R$ zur reflexiv-transitiven Hülle $R_{trans}$ von $R$ (also nochmal neu zeichnen als reflexiv-transitive Hülle).
    \item $R_{trans}$ ist eine Ordnungsrelation. Was sind die minimalen, maximalen, kleinsten, gr"o{\ss}ten Elemente (sofern existent) davon? Begründen Sie Ihre Antwort.

\end{enumerate}


\aufgabetitel{$6$}{"Aquivalenzrelationen und Partitionen}\\
Sei $M=\{1,2,3,4,5\}$, $A=\{1,2,3\}$, also $A\subseteq M$ und die Partition $P$ auf $M$ mit $P=\{A, \{4\}, \{5\}\}$ gegeben.
\begin{enumerate}[(i)]
    \item Warum ist $P$ eine Partition auf $M$? 
    \item Eine Partition induziert eine "Aquivalenzrelation. Definieren Sie die zu $P$ passend "Aquivalenzrelation.
    \item Argumentieren Sie warum Ihre Relation eine "Aquivalenzrelation ist.
    \item Geben Sie die "Aquivalenzklasse $[1]$ an.
\end{enumerate}



\aufgabetitel{$4$}{Indexmengen und Beweis}\\
Es sei $A_i=\{n\in\mathbb{N}\mid n<i\}$.
\begin{enumerate}[(i)]
    \item Bestimmen Sie $A_4$.
    \item Zeigen Sie, dass $\bigcap_{i\in\mathbb{N}}A_i=\emptyset$.
\end{enumerate}

\aufgabetitel{$5$}{Funktionen und Relationen}\\
Sei $f:X\to Y$ eine Funktion und $\sim$ eine Relation "uber $X$ mit $x\sim y:\Leftrightarrow f(x)\leq f(y)$.
\begin{enumerate}[(i)]
    \item Stellen Sie die Funktion $f$ und die Relation $\sim$ f"ur $X=\{1,2,3,4,5\}$, $Y=\{1,2,3\}$, $1\mapsto 1, 2\mapsto 1, 3\mapsto 2, 4\mapsto 2, 5\mapsto 3$ dar.
    \item Zeigen Sie: $\sim$ ist eine Ordnungsrelation genau dann, wenn $f$ injektiv ist.
\end{enumerate}
\hspace{5cm}
\newpage
\underline{\textbf{Formelsammlung}}\\
Hier eine kleine Formelsammlung. Sie ist nicht vollständig, enthält aber alle wichtigen Statements / Definitionen, die man brauchen könnte.
\begin{enumerate}
\item Aussagen- und Pr"adikatenlogik
\begin{enumerate}
    \item Distributivgesetz: $A\wedge (B \vee C) \Leftrightarrow (A\wedge B) \vee (A\wedge C)$
    \item Distributivgesetz: $A\vee (B \wedge C) \Leftrightarrow (A\vee B) \wedge (A\vee C)$
    \item DeMorgan: $\neg(A\wedge B) \Leftrightarrow \neg A \vee \neg B$
    \item DeMorgan: $\neg(A\vee B) \Leftrightarrow \neg A \wedge \neg B$
    \item Idempotenz: $A\wedge A \Leftrightarrow A$
    \item Idempotenz: $A\vee A \Leftrightarrow A$
    \item $A\wedge \neg A \Leftrightarrow \bot$
    \item $A\vee \neg A \Leftrightarrow \top$
    \item $\neg\neg A \Leftrightarrow A$
    \item $\neg\forall x\in M: A(x) \Leftrightarrow \exists x\in M: \neg A(x)$
    \item $\neg\exists x\in M: A(x) \Leftrightarrow \forall x\in M: \neg A(x)$
\end{enumerate}
\item Mengen
\begin{enumerate}
    \item Teilmenge: $A\subseteq B \Leftrightarrow \forall x\in A: x\in B$
    \item Potenzmenge: $\mathscr{P}(A) = \{B\mid B\subseteq A\}$
    \item Vereinigung: $A\cup B = \{x\mid x\in A \vee x\in B\}$
    \item Schnittmenge: $A\cap B = \{x\mid x\in A \wedge x\in B\}$
    \item Differenzmenge: $A\setminus B = \{x\mid x\in A \wedge x\notin B\}$
    \item Distributivgesetz: $A\cap (B \cup C) \Leftrightarrow (A\cap B) \cup (A\cap C)$
    \item Distributivgesetz: $A\cup (B \cap C) \Leftrightarrow (A\cup B) \cap (A\cup C)$
    \item DeMorgan: $A\setminus (B\cup C) \Leftrightarrow (A\setminus B) \cap (A\setminus C)$
    \item DeMorgan: $A\setminus (B\cap C) \Leftrightarrow (A\setminus B) \cup (A\setminus C)$
    \item Es ist $\bigcup_{i\in I}A_i = \{x\mid \exists i\in I: x\in A_i\}$. 
    \item Es ist $\bigcap_{i\in I}A_i = \{x\mid\forall i\in I: x\in A_i\}$.
\end{enumerate}
\item Relationen
\begin{enumerate}
    \item F"ur Mengen $M,N$ ist $R\subseteq M\times N$ eine Relation von $M$ nach $N$.
    \item $R\subseteq M\times M$ ist reflexiv, wenn $\forall x\in M: (x,x)\in R$.
    \item $R\subseteq M\times M$ ist symmetrisch, wenn $\forall x,y\in M: (x,y)\in R \Rightarrow (y,x)\in R$.
    \item $R\subseteq M\times M$ ist antisymmetrisch, wenn $\forall x,y\in M: (x,y)\in R \wedge (y,x)\in R \Rightarrow x=y$.
    \item $R\subseteq M\times M$ ist transitiv, wenn $\forall x,y,z\in M: (x,y)\in R \wedge (y,z)\in R \Rightarrow (x,z)\in R$.
    \item $R\subseteq M\times M$ ist eine "Aquivalenzrelation, wenn $R$ reflexiv, symmetrisch und transitiv ist.
    \item $R\subseteq M\times M$ ist eine Ordnungsrelation, wenn $R$ reflexiv, antisymmetrisch und transitiv ist.
    \item F"ur eine "Aquivalenzrelation $\sim$ auf $M$ ist $[x]=\{y\in M\mid x\sim y\}$ die "Aquivalenzklasse von $x$, $M/\sim = \{[x]\mid x\in M\}$ die Menge der "Aquivalenzklassen oder Quotientenmenge von $M$ modulo $\sim$. Die Menge der "Aquivalenzklassen ist eine Partition von $M$. Umgekehrt induziert jede Partition eine "Aquivalenzrelation.
    \item F"ur eine Ordnungsrelation $\leq$ auf $M$ und $X\subseteq M$ ist $g$ ein kleinstes Element von $X$, wenn $\forall x\in X: g\leq x$, $g$ ein minimales Element von $X$, wenn $\forall g'\in X: g'\leq g\Rightarrow g=g'$, maximale und gr"o{\ss}te Elemente analog.
\end{enumerate}
\item Funktionen
\begin{enumerate}
    \item Eine Funktion $f:X\to Y$ ist eine Relation (also $f\subseteq X\times Y$), die jedem Element aus der Definitionsmenge $X$ genau ein Element aus der Zielmenge $Y$ zuordnet.
    \item $f$ ist injektiv, wenn $\forall x_1,x_2\in X: f(x_1)=f(x_2) \Rightarrow x_1=x_2$.
    \item $f$ ist surjektiv, wenn $\forall y\in Y: \exists x\in X: f(x)=y$.
    \item $f$ ist bijektiv, wenn $f$ injektiv und surjektiv ist.
    \item Die Umkehrfunktion $f^{-1}$ ist definiert $f^{-1}(Y)=\{x\in X\mid\exists y\in Y: y=f(x)\}$
\end{enumerate}
\end{enumerate}
\end{document}